% Part: computability
% Chapter: computability-theory
% Section: total

\documentclass[../../../include/open-logic-section]{subfiles}

\begin{document}

\olfileid{cmp}{thy}{tot}
\olsection{تام‌بودن تصمیم‌ناپذیر است}

نمونهٔ دیگری از کاربرد قضیهٔ $s$-$m$-$n$ برای نشان‌دادن
محاسبه‌ناپذیری چیزی را بررسی کنیم. بگذارید $\fn{Tot}$ مجموعهٔ
نمایه‌های توابع محاسبه‌پذیر تام باشد، یعنی
\[
\fn{Tot} = \Setabs{x}{\text{به ازای هر $y$، $\cfind{x}(y)\fdefined$}}.
\]

\begin{prop}
\ollabel{prop:total}
$\fn{Tot}$ محاسبه‌پذیر نیست.
\end{prop}

\begin{proof}
برای دیدن محاسبه‌ناپذیریِ $\fn{Tot}$ کافی است نشان دهیم $K$ به آن
کاهش‌پذیر است. $h(x,y)$ را چنین تعریف کنید:
\[
h(x,y) \simeq
\begin{cases}
0 & \text{اگر $x \in K$} \\
\fundefined & \text{در غیر این صورت.}
\end{cases}
\]
توجه کنید که $h(x,y)$ اصلاً به~$y$ وابسته نیست. دیدن اینکه $h$
محاسبه‌پذیر جزئی است نباید دشوار باشد: روی ورودی‌های $x,y$ نخست تابع
$\cfind{x}$ را روی ورودی~$x$ شبیه‌سازی می‌کنیم؛ اگر این محاسبه متوقف
شود، $h(x,y)$ مقدار~$0$ را بیرون می‌دهد و متوقف می‌شود. پس
$h(x,y)$ همان $\Zero(\umin{s}{T(x,x,s)})$ است، که در آن $\Zero$
تابع ثابت صفر است.

با استفاده از قضیهٔ $s$-$m$-$n$، تابع بازگشتی اولیه‌ای چون~$k(x)$
وجود دارد که به ازای هر $x$ و~$y$،
\[
\cfind{k(x)}(y) \simeq
\begin{cases}
0 & \text{اگر $x \in K$} \\
\fundefined & \text{در غیر این صورت.}
\end{cases}
\]
پس اگر $x\in K$، تابع $\cfind{k(x)}$ تام است و در غیر این صورت
هیچ‌جا تعریف نشده است. بنابراین $k$ کاهشی از $K$ به~$\fn{Tot}$ است.
\end{proof}

\begin{digress}
در واقع $\fn{Tot}$ حتی محاسبه‌پذیرِ برشمردنی هم نیست---پیچیدگی آن در
«سلسله‌مراتب حسابی» بالاتر قرار دارد. اما اینجا به این تقویت
نمی‌پردازیم.
\end{digress}

\end{document}
